\documentclass[11pt,a4paper]{article}
\usepackage{fontspec,amsmath,amssymb,amsthm,microtype}
\setmainfont{lmroman10-regular.otf}[BoldFont=lmroman10-bold.otf,ItalicFont=lmroman10-italic.otf,BoldItalicFont=lmroman10-bolditalic.otf]
\setmonofont{lmmono10-regular.otf}[Scale=MatchLowercase]
\usepackage[margin=25mm]{geometry}
\usepackage[unicode,hidelinks]{hyperref}
\setlength{\emergencystretch}{3em}
\setlength{\parskip}{0.4em}
\setlength{\parindent}{0pt}
\allowdisplaybreaks[2]
\raggedbottom
\title{\texorpdfstring{The shared exponential proof at particle temperature}{The shared exponential proof at particle temperature}}
\author{LangevinLean mathematical exposition}
\date{}
\begin{document}\maketitle
\noindent Source module: \texttt{KIPLMC/\allowbreak{}Reference\allowbreak{}Bound.\allowbreak{}lean}. The accompanying HTML pairs this mathematical exposition with the exact Lean statements and records the scope of the formalization.


The first particle algorithm uses exactly the shared exponential step. Its analytical comparison should therefore reuse the same proof. This module applies the shared L\ensuremath{{}^{2}} kinetic theorem to the averaged lifted potential, uses the actual thermal Gaussian radius, and projects the resulting transport estimate to the MMLE parameter. The result has all three constants in the supplied KIPLMC1 theorem, with the initial supplied coupling cost. Its concrete integral-segment and reference-marginal hypotheses are explicit: construction of a canonical stationary reference and identification of the independent Gaussian chain are still needed for the unconditional sampler theorem.


\section{\texorpdfstring{Lift the actual physical particle force once}{Lift the actual physical particle force once}}


The averaged lifted potential \(\bar U_N=V_N/N\) has gradient \(F_N\), the physical-time force. The earlier geometry proves its original smoothness constant \(L\) and strong convexity constant \(m\), independently of particle count. \texttt{lp\allowbreak{}Force} applies the existing pointwise \texttt{ULA.\allowbreak{}lp\allowbreak{}Drift} construction to this gradient. It is therefore the actual particle force on almost-everywhere representatives, with a well-defined square-integrable output. No new exponential update or substitute gradient algorithm is introduced.



{\raggedright\noindent Lean declarations: \texttt{lp\allowbreak{}Force}.\par}

\section{\texorpdfstring{Inherit the local comparison, recurrence, and sharp equilibrium term}{Inherit the local comparison, recurrence, and sharp equilibrium term}}


Let \(S_n\) be numerical phase states and \(R_n\) reference phase states in L\ensuremath{{}^{2}}. On each interval, the auxiliary exact path starts at \(S_n\), the reference exact path starts at \(R_n\) and ends at \(R_{n+1}\), and their integral equations have the same additive input. The numerical endpoint is the actual shared exponential step with velocity noise \(B_h\) and position noise \(\int_0^hB_sds\). Its mild velocity and kinematic identities are explicit inputs. The reference velocity marginal is the thermal Gaussian and its position marginal is the normalized lifted Gibbs target.



The thermal moment theorem supplies \(C=\sqrt{(d_\theta+Nd_x)/N}\). Applying the shared kinetic theorem gives the full-state transport estimate with initial term \(\sqrt2(1-3mh/(4\gamma))^n\|S_0^X-R_0^X\|_{L^2}\) and bias \(Lh\sqrt2C/m\), under \(\gamma^2\ge m+L\) and \(4\gamma Lh\le m\). Every local-error and recurrence step is inherited, not assumed as a transport inequality.



The particular L\ensuremath{{}^{2}} pair \((S_n^X,R_n^X)\) supplies a finite transport coupling, so real-valued conversion cannot conceal infinity. The parameter projection has operator norm at most one. The earlier sharp Gibbs concentration theorem contributes \(\sqrt{d_\theta/(Nm)}\). Lean invokes the proved transport-to-projected-RMS triangle bound, proves the projected moment is finite, rewrites it as the square root of the actual parameter second moment, and adds the three terms. The resulting theorem remains explicitly a reference comparison; it does not assert that the unresolved reference construction has already been discharged.



{\raggedright\noindent Lean declarations: \texttt{kiplmc1\_\allowbreak{}reference\_\allowbreak{}parameter\_\allowbreak{}bound}.\par}

\end{document}
